\documentclass[11pt,a4paper]{article}
\usepackage[utf8]{inputenc}
\usepackage[margin=2.2cm]{geometry}
\usepackage{graphicx,booktabs,longtable,array,listings,xcolor}
\usepackage[hidelinks]{hyperref}
\newcommand{\code}[1]{\path{#1}}
\newcommand{\todo}[1]{\textcolor{red}{[#1]}}
% Include an image if the file exists, otherwise print a red note instead of failing.
\newcommand{\maybeimg}[2]{\IfFileExists{#1}{\includegraphics[width=#2]{#1}}{\todo{missing image: \detokenize{#1}}}}
\lstset{basicstyle=\ttfamily\footnotesize,frame=single,breaklines=true,columns=fullflexible}
\setlength{\parindent}{0pt}\setlength{\parskip}{4pt}

\begin{document}

\begin{center}
{\LARGE\bfseries 23CSE455 Design Patterns\\Case Study Report}\\[6pt]
\end{center}

\begin{tabular}{@{}ll@{}}
\textbf{Team:} & B14\\
\textbf{Case study:} & Design Pattern Identification and Evolution in a Next.js-Based\\
 & Digital News and Publishing Application\\
\textbf{Application:} & Next.js Blog CMS\\
\textbf{Date:} & 2026-10-28 (Part 3 records corrected)\\
\textbf{Submission repository:} & 23CSE455\_DP\_CaseStudy
\end{tabular}

\subsection*{Students}
\begin{tabular}{@{}lll@{}}
\toprule
Roll number & Student name & Batch\\
\midrule
AM.SC.U4AIE23121 & Keziah Suzaine George & AIE\\
AM.SC.U4AIE23139 & Nivedita Ramesh & AIE\\
AM.SC.U4CSE23360 & Soorya Krishna & CSE\\
AM.SC.U4CSE23309 & B R Govind & CSE\\
\bottomrule
\end{tabular}

%------------------------------------------------------------
\section{Application and architecture}

\textbf{Original source.}
Repository: \url{https://github.com/thepracticalengineer/nextjs-blog-cms.git}.
Baseline commit: \code{90ec6565877053ee0e93e9f09dcc7c7213aa4832}
(merge of pull request \#102, 8 October 2026). Access date: 2026-10-10. License: MIT.

\textbf{Application.} A full-stack blog CMS with a TipTap editor, Supabase
authentication including MFA, Admin and Author roles, a draft and publication
workflow, comments, tags and categories, public author profiles, newsletter
subscriptions, an AI writing assistant (Claude, Gemini, OpenAI) and a developer REST API.

\textbf{Environment.}
\begin{tabular}{@{}ll@{}}
Node.js & 24.14.0 (repository requests 24.21.0; mismatch noted, nothing failed)\\
Package manager & pnpm 10.29.3\\
Framework & Next.js 16.3.8 (App Router), React 19.3.0, TypeScript 6.0.3\\
Database and auth & Supabase (supabase-js 2.117.2, @supabase/ssr 0.12.7)\\
Other & Zustand 5.0.15, Zod 4.6.5, Vitest 5.0.3, Playwright 1.63.0\\
\end{tabular}

\textbf{Scope.} Inspected: \code{app/}, \code{components/}, \code{features/},
\code{lib/} and \code{proxy.ts}. Excluded: \code{node_modules}, \code{.next},
\code{.git}, \code{.husky}, \code{.github}, \code{.agents}, \code{.claude},
\code{docs/}, \code{public/}, \code{database/}, \code{supabase/migrations/},
\code{e2e/}, \code{__tests__/}, lockfiles and generated files. Reason: not
application logic or not relevant to pattern discovery. Unit tests were used only
as supporting evidence. The source copy in \code{code/} equals the baseline commit
minus \code{.husky}, \code{.github}, \code{tsconfig.tsbuildinfo} and
\code{package-lock.json}.

\textbf{Architecture.}
\begin{figure}[htbp]
\centering
\maybeimg{Part1/architecture.png}{\linewidth}
\caption{Architecture of the application at the baseline commit.}
\end{figure}

The browser reaches public pages, dashboard pages, authentication pages and API
routes. Dashboard and authentication paths pass through \code{proxy.ts} first; its
\code{matcher} covers only \code{/dashboard}, \code{/login}, \code{/register},
\code{/mfa}, \code{/forgot-password} and \code{/reset-password}, so public pages do
not. Pages and routes call domain modules in \code{features/}, which use the
Supabase client creators in \code{lib/supabase}, the permission table in
\code{lib/permissions}, the LLM service and the notification modules.

\textbf{Request flow (admin opens the Users page).} The browser requests
\code{/dashboard/admin/users}. Next.js first runs \code{proxy.ts}, which refreshes
the Supabase session with \code{supabase.auth.getUser()} (line 34), redirects
unauthenticated users to \code{/login}, sends users with an enrolled but unverified
second factor to \code{/mfa}, and, for paths under \code{/dashboard/admin}, reads
the user's role from the \code{profiles} table and redirects non-admins to
\code{/dashboard}. When the request passes, the server component \code{UsersPage}
(\code{app/(dashboard)/dashboard/admin/users/page.tsx}) runs. It calls
\code{requirePermission('users:read')} (line 9), loads data with
\code{getAllUsers()} from \code{features/users/queries} (line 11) and renders
\code{UserTable} (line 21). In the browser, \code{AuthProvider} keeps session state
current by observing \code{supabase.auth.onAuthStateChange}.

\textbf{Run instructions.} See \code{README.md}. Results at the baseline:

\begin{tabular}{@{}p{5.6cm}p{9.6cm}@{}}
\toprule
Command & Result\\
\midrule
\code{pnpm install --frozen-lockfile} & Succeeded, 933 packages\\
\code{pnpm run test:run} & 927 of 927 passed in 83 files on default settings (runs of 71.43 s and 80.10 s). Log: \code{test_result/test_result_vitest.log}\\
\code{pnpm exec vitest run --testTimeout=30000 --maxWorkers=2} & 83 test files passed, 927 of 927 tests passed (222.58 s). Screenshot: \code{test_result/vitest_summary.png}\\
\code{pnpm run typecheck} & Passed. Log: \code{test_result/typecheck.log}; screenshot: \code{test_result/typecheck.png}\\
\code{pnpm run dev}, \code{pnpm run test:e2e} & Not run: Docker is not installed, so local Supabase was unavailable\\
\bottomrule
\end{tabular}

Two further default-setting runs under heavy machine load failed with 6 and 12
failures (130 s and 188 s). Every failure was a 5000 ms timeout in a jsdom UI test,
or a follow-on assertion in the same file after a timeout. The repository sits in a
OneDrive-synced folder. With a 30 s timeout and 2 workers the same code passed
927 of 927, so the failures were timing-related and not application defects.

\begin{figure}[htbp]
\centering
\maybeimg{test_result/vitest_summary.png}{0.9\linewidth}\\[4pt]
\maybeimg{test_result/typecheck.png}{0.9\linewidth}
\caption{Passing unit-test run and typecheck at the baseline commit.}
\end{figure}

%------------------------------------------------------------
\section{Part 1: Pattern discovery}

\textbf{Method.} Candidates were found by reading the file listing and then
reading the likely files in full with line numbers (saved in
\code{Part1/evidence/}). Each candidate was checked against three questions: what
problem it solves here, whether every role maps to a real symbol, and whether the
participants interact as the pattern defines. A class name, a file name or an
LLM statement alone was not accepted as evidence. Combined or absent roles are
stated in the records; no classes were invented.

\subsection{Totals at the baseline}
\begin{tabular}{@{}p{5.5cm}p{9.7cm}@{}}
\toprule
Confirmed pattern types & 3 (Strategy, Observer, Role-Based Access Control)\\
Confirmed pattern instances & 3\\
Concrete participant classes & 7, all functions or components rather than classes: \code{streamClaude}, \code{streamOpenAI}, \code{streamGemini}; \code{AuthProvider}; \code{can}, \code{canAccess}, \code{canAny}\\
Unique participant symbols & 15 application symbols, plus 1 from a dependency (\code{supabase.auth.onAuthStateChange})\\
Rejected candidates & 7\\
Uncertain candidates & 1\\
Observation scope & Files listed in \code{Part1/rejected_uncertain.md}, at the baseline commit\\
\bottomrule
\end{tabular}

\subsection{Pattern inventory}
\begin{tabular}{@{}llcc@{}}
\toprule
Pattern type & Instance name & Instances & Concrete classes\\
\midrule
Strategy & LLM Provider Streaming Strategy & 1 & 3\\
Observer & Auth State Observer & 1 & 1\\
RBAC & RBAC Permission Table & 1 & 3\\
\bottomrule
\end{tabular}

\subsection{Records for each pattern instance}

%---------------- Strategy
\subsubsection*{Instance 1: LLM Provider Streaming Strategy}
\textbf{Pattern type and category:} Strategy, GoF behavioral.
\textbf{Purpose and origin:} streams chat replies from Claude, OpenAI or Gemini
through one call; application code.

\begin{longtable}{@{}p{3.2cm}p{5.8cm}p{6.2cm}@{}}
\toprule
Role & Symbol & File and lines\\
\midrule
\endhead
Strategy & Not separately declared (implicit signature \code{(params: SendMessageParams) => AsyncGenerator<string>}) & \code{features/ai-assistant/llmService.ts}:10--16 (params type)\\
ConcreteStrategy & \code{streamClaude} & same file:69--90\\
ConcreteStrategy & \code{streamOpenAI} & same file:94--118\\
ConcreteStrategy & \code{streamGemini} & same file:129--146\\
Context & \code{sendMessage} & same file:150--155\\
\bottomrule
\end{longtable}

\textbf{Evidence.} The three functions share one parameter type and return type.
The context selects one at run time from \code{params.provider} and delegates:
\begin{lstlisting}
export function sendMessage(params: SendMessageParams): AsyncGenerator<string> {
  if (params.provider === 'claude') return streamClaude(params)
  if (params.provider === 'openai') return streamOpenAI(params)
  if (params.provider === 'gemini') return streamGemini(params)
  throw new Error(`Unsupported provider: ${params.provider}`)
}
\end{lstlisting}
\textbf{Limitation.} There is no declared interface, and selection is a hard-coded
\code{if} chain rather than injection or a registry.

\begin{figure}[htbp]\centering
\maybeimg{Part1/uml_strategy.png}{0.85\linewidth}
\caption{UML for the Strategy instance (source: \texttt{Part1/uml\_strategy.mmd}).}
\end{figure}

\textbf{Verification:} Confirmed. Reviewers: team B14. Method: source reading of
the listed ranges; \code{__tests__/ai-assistant/llmService.test.ts} (20 tests)
passes in the baseline run but was not the sole evidence.
Alternative considered: the repeated provider branches in
\code{generateBlogPost}, \code{generateBlogPostHeadless}, \code{generateChatTitle}
and \code{validateProviderKey} (lines 157--386) do not apply Strategy; see the
rejected list.

%---------------- Observer
\subsubsection*{Instance 2: Auth State Observer}
\textbf{Pattern type and category:} Observer, GoF behavioral.
\textbf{Purpose and origin:} keeps UI session and profile state in sync with login,
logout and token refresh; application code using the Supabase dependency.

\begin{longtable}{@{}p{3.2cm}p{5.8cm}p{6.2cm}@{}}
\toprule
Role & Symbol & File and lines\\
\midrule
\endhead
Subject & \code{supabase.auth} (dependency, not application code) & \code{features/auth/context/AuthProvider.tsx}:37--39\\
ConcreteObserver & \code{AuthProvider}, via the callback given to \code{onAuthStateChange} & same file:9--59 (callback 39--48)\\
Detach & \code{subscription.unsubscribe()} & same file:50\\
State holder & \code{AuthContext}, \code{AuthContext.Provider} & \code{AuthContext.tsx}:14--19; \code{AuthProvider.tsx}:54--58\\
Consumers & \code{useAuth}; \code{useSession} & \code{AuthContext.tsx}:21--23; \code{features/auth/hooks/useSession.ts}:5--8\\
\bottomrule
\end{longtable}

\textbf{Evidence.} The component subscribes, updates state when the callback
fires, and unsubscribes in the effect cleanup:
\begin{lstlisting}
const { data: { subscription } } = supabase.auth.onAuthStateChange(
  async (_event, session) => { setSession(session); setUser(session?.user ?? null); ... })
return () => subscription.unsubscribe()
\end{lstlisting}
The subject belongs to \code{@supabase/supabase-js}, so no subject class is
invented. React Context is the delivery mechanism to consumers and is not counted
as a second instance.

\begin{figure}[htbp]\centering
\maybeimg{Part1/uml_observer.png}{0.85\linewidth}
\caption{UML for the Observer instance (source: \texttt{Part1/uml\_observer.mmd}).}
\end{figure}

\textbf{Verification:} Confirmed. Reviewers: team B14. Method: source reading of
\code{AuthProvider.tsx}, \code{AuthContext.tsx} and \code{useSession.ts}.

%---------------- RBAC
\subsubsection*{Instance 3: RBAC Permission Table}
\textbf{Pattern type and category:} Role-Based Access Control, architecture and
security pattern (not a GoF pattern). \textbf{Purpose and origin:} centralizes
which role may perform which action; application code.

\begin{longtable}{@{}p{3.2cm}p{5.8cm}p{6.2cm}@{}}
\toprule
Role & Symbol & File and lines\\
\midrule
\endhead
Permission table & \code{rolePermissions} & \code{lib/permissions/index.ts}:3--30\\
Check & \code{can} & same file:32--35\\
All-of check & \code{canAccess} & same file:37--40\\
Any-of check & \code{canAny} & same file:42--45\\
Types & \code{Role}; \code{Permission} & \code{lib/permissions/types.ts}:1; 3--18\\
\bottomrule
\end{longtable}

\textbf{Evidence.} \code{rolePermissions} maps each of the roles \code{admin} and
\code{author} to a list of permissions. \code{can} looks up the table, and
\code{canAccess} and \code{canAny} apply \code{every} and \code{some} over
\code{can}:
\begin{lstlisting}
export function can(role: Role | null | undefined, permission: Permission): boolean {
  if (!role) return false
  return rolePermissions[role]?.includes(permission) ?? false
}
\end{lstlisting}

\begin{figure}[htbp]\centering
\maybeimg{Part1/uml_rbac.png}{0.8\linewidth}
\caption{UML for the RBAC instance (source: \texttt{Part1/uml\_rbac.mmd}).}
\end{figure}

\textbf{Verification:} Confirmed as an architecture pattern. Reviewers: team B14.
Method: source reading; \code{__tests__/lib/permissions.test.ts} (29 tests) passes
in the baseline run.

\subsection{Rejected and uncertain candidates}
\textbf{Uncertain (1).} \code{proxy.ts} guard pipeline (lines 4--117): sequential
guard checks in one function. It resembles Chain of Responsibility, but there is no
handler interface, no successor links and no separate handlers, so roles cannot be
mapped without inventing classes. Classified as a Next.js framework idiom.

\textbf{Rejected (7).}
\begin{longtable}{@{}cp{5cm}p{9.5cm}@{}}
\toprule
\# & Candidate & Reason\\
\midrule
\endhead
R1 & Supabase client creators as Factory Method & Four standalone functions chosen by import path; no creator hierarchy or polymorphic selection; \code{createStaticClient} can return null, so they share no contract\\
R2 & User-confirmed notifications as Observer & \code{route.ts}:67--74 calls two functions in \code{Promise.all}, each with its own catch; no subscriber list or notifier interface\\
R3 & \code{sendNewsletterEmail} & Single function wrapping Resend; one participant\\
R4 & Publication validation as State or Specification & \code{status} is a plain string; one function with sequential checks, no rule objects\\
R5 & Repeated provider \code{if} chains & Strategy is not applied; branching duplicated (Part 2 refactoring target)\\
R6 & Helper functions & \code{buildSystemPrompt}, \code{resolveGeminiModel}, \code{isRateLimitError}, \code{escapeHtml}, \code{sanitizeSubject}: plain utilities\\
R7 & React Context as a separate Observer & Folded into the Auth State Observer to avoid double counting\\
\bottomrule
\end{longtable}

%------------------------------------------------------------
\section{Part 2: Evolution and refactoring}

\textbf{Scope and traceability.} Part 2 evolves the application from the baseline
source snapshot \code{de5bc25} in three requested areas: related-article
recommendations, content moderation, and premium-content access control. The
implementation is maintained in the application fork
\url{https://github.com/nivedita-rmsh/nextjs-blog-cms}, branch
\code{part2/content-moderation-paywall}. The latest inspected commit is
\code{7ecfa65}. The report's baseline copy in \code{B14/code/} remains the
Part 1 snapshot; it is not a copy of the evolved application branch. The Part 1
pattern inventory is tied to commit
\code{90ec6565877053ee0e93e9f09dcc7c7213aa4832}, while the Part 2 feature work
starts from the later application commit \code{de5bc25}. These snapshots must
not be conflated when describing the starting state.

\begin{tabular}{@{}p{1.5cm}p{3.2cm}p{2.4cm}p{7.3cm}@{}}
\toprule
Change & Requirement & Evidence commit & Main implementation areas\\
\midrule
CH01 & Article recommendations & \code{68cd127} & \code{features/posts/recommendations.ts}; \code{features/posts/components/RelatedPosts.tsx}; public article page\\
CH02 & Content moderation & \code{376bc05}--\code{33ca044} & Moderation state, admin review page and actions, public queries, REST publication paths, Supabase migration\\
CH03 & Premium-content access & \code{a7d6990}--\code{7ecfa65} & Premium entitlement helper, editor toggle, protected body storage, public article page and Supabase migration\\
\bottomrule
\end{tabular}

The changes are cumulative on the feature branch. The short commit IDs above
identify the first or final relevant commits; the complete history is available
in the application repository. Code paths and line ranges below refer to that
branch, rather than the Part 1 baseline.

\subsection*{CH01: Related-article recommendations}
\textbf{Requirement.} Show a small set of useful related articles on an article
detail page, based on existing metadata, without recommending the current
article, drafts, unapproved posts, or unrelated content.

\textbf{Implementation.} \code{rankRelatedPosts()} in
\code{features/posts/recommendations.ts} (lines 7--49) ranks candidates by
shared tags and category. Each shared tag contributes two points and a shared
category contributes one point; candidates with no positive score are removed.
Ties are resolved by publication recency, and the result is limited to three
posts by default. The function excludes the current article, drafts and unrelated candidates. The later CH02 moderation change adds the additional approval-state filter to recommendation candidates in the cumulative branch. The async server component
\code{RelatedPosts} in \code{features/posts/components/RelatedPosts.tsx}
(lines 6--23) loads published posts, applies the ranker and renders the results
when any remain. The public detail page imports and renders the component at
\code{app/(public)/blog/[slug]/page.tsx:14,143}.

\textbf{Pattern effect.} No GoF pattern is claimed as newly introduced. The
ranker is a deterministic domain helper, not a Strategy instance: there is no
family of interchangeable ranking algorithms selected through a common
contract. The existing LLM Strategy, Auth State Observer and RBAC Permission
Table are preserved. The added query/component collaboration is recorded as a
feature-level change, not counted as another design-pattern instance.

\textbf{Verification.} The focused file \code{__tests__/features/posts/recommendations.test.ts}
contains four cases covering tag-versus-category ranking, exclusion of the
current article and drafts, result limiting/relevance, and empty results. The
recorded focused run showed all 4 recommendation tests passing. During the same
pre-commit attempt, a separate author-byline test file reported two failures;
those failures were not established as pre-existing, so the overall hook result
must not be described as fully passing. The commit was created with hooks
bypassed after the hook failure. A full passing post-change suite has not been
recorded for this snapshot.

\subsection*{CH02: Content moderation before publication}
\textbf{Requirement.} Require administrator review before new publication
requests become public, apply that rule to the normal editor and supported REST
publication paths, and keep pending or rejected content out of public views.

\textbf{Implementation.} The \code{moderation_status} column supports
\code{not_submitted}, \code{pending}, \code{approved} and \code{rejected}.
The migration \code{supabase/migrations/20261012100000_content_moderation.sql}
(lines 1--117) backfills existing published posts as approved, defines the
public-read policies, and updates the atomic save function so a publication
request is stored as a draft pending review. Edits to a published article return
it to the queue. Public post queries, category/tag listings, sitemap entries,
recommendation candidates and the relevant REST routes filter or enforce the
approval state; examples include \code{features/posts/queries.ts} (lines 30,
48, 96, 117 and 148), \code{app/api/posts/create/route.ts} (lines 29--33,
75--91 and 116--119), and \code{app/api/posts/[id]/route.ts} (lines 138--140
and 201--208).

Administrators see \code{/dashboard/admin/moderation}, implemented in
\code{app/(dashboard)/dashboard/admin/moderation/page.tsx} (lines 8--39), with
the review controls in \code{components/dashboard/ModerationQueue.tsx}
(lines 20--53). The server actions in \code{features/posts/moderation-actions.ts}
(lines 12--39 and 42--62) perform admin checks, revalidate publication fields,
compare the row's update timestamp to avoid acting on stale content, and then
approve/publish or reject/return the post to draft. The editor shows a pending
moderation message instead of claiming publication succeeded.

\textbf{Pattern effect.} The original RBAC Permission Table remains present; its
three baseline participants are still the same. The moderation feature adds
admin-role checks and database row-level security (RLS) at new enforcement
points. These controls are not a new GoF pattern, and the new explicit admin
checks are not presented as a complete centralization of authorization through
\code{rolePermissions}. The existing Strategy and Observer instances are
preserved.

\textbf{Verification status.} Five unit cases were added in
\code{__tests__/features/posts/moderation.test.ts} for state transitions and
public visibility, and existing post-action and REST tests were adjusted.
These focused tests were added but their run result had not been captured at the
time of this report update. The migration has not yet been applied and verified
against a running Supabase database. TypeScript checking passed on the cumulative
feature branch, but this is not a substitute for running the tests or the
migration.

\subsection*{CH03: Premium-content access control}
\textbf{Requirement.} Allow an author to mark an article premium, show a public
preview to readers without access, and disclose the full body only to the author,
an administrator or a user with an active premium entitlement. Newsletter
subscription is not treated as proof of paid access.

\textbf{Implementation.} The editor schema and \code{PostEditor} add an
\code{is_premium} flag and checkbox (\code{components/dashboard/PostEditor.tsx}:
38, 122 and 482--485). The migration
\code{supabase/migrations/20261012103000_premium_content_access.sql}
(lines 1--170) adds the flag to \code{posts}, creates
\code{premium_post_contents} for full premium bodies and
\code{premium_subscriptions} for access entitlements, and adds RLS policies for
approved premium content and active entitlements. The atomic persistence
function stores a premium post's full body in the protected table and leaves
\code{posts.content} empty for that post; when a post is changed back to free,
the full body is restored to the normal field and the premium-body row is
removed.

\code{hasPremiumAccess()} and \code{getPremiumPostBody()} in
\code{features/premium/access.ts} (lines 5--28) check sign-in, admin/author
access, entitlement status and optional expiry before retrieving the protected
body. The public article page (\code{app/(public)/blog/[slug]/page.tsx}, lines
58--60 and 119--128) renders the full content only when access is confirmed;
otherwise it renders an excerpt and sign-in call to action. The REST detail
endpoint remains API-key protected and scopes a post read to the API-key owner's
posts; it is used for author operations, not anonymous public paywall access.

\textbf{Scope limitation.} This change implements entitlement-based access
control, not a payment or billing system. An active row in
\code{premium_subscriptions} must be provisioned through an administrative or
separate trusted process; the application does not process payments, create
paid subscriptions, or handle renewals.

\textbf{Pattern effect.} The existing Strategy, Observer and RBAC instances
remain present. Premium entitlement checks and protected-content persistence
are additional security/domain rules, not a demonstrated GoF pattern. No new
pattern is counted based solely on the helper's name or on the presence of a
subscription table.

\textbf{Verification status.} Six test cases were added in
\code{__tests__/features/premium/access.test.ts} for anonymous denial, active and
expired entitlements, revoked entitlements, administrator access, author access
to a protected body and denial before protected content is queried. The test file
was added but its run result had not been captured at the time of this report
update. The cumulative branch's TypeScript command
\code{pnpm exec tsc --noEmit} completed with exit code 0. The new Supabase
migrations and premium-content flows still require database/runtime verification;
no full-suite passing result is claimed for this branch.

\subsection*{Pattern evolution summary}
The baseline inventory contains three confirmed instances: LLM Provider
Streaming Strategy, Auth State Observer and RBAC Permission Table. Source
inspection of the feature branch does not establish a fourth GoF instance, so
these counts remain three across the snapshots below. The table describes the
three confirmed Part 1 instances only; it does not count every feature module.

\begin{tabular}{@{}p{3.3cm}p{2.4cm}ccc p{3.0cm}@{}}
\toprule
Snapshot & Commit & Strategy & Observer & RBAC & Total confirmed instances\\
\midrule
Baseline & \code{de5bc25} & 1 & 1 & 1 & 3\\
After CH01 recommendations & \code{68cd127} & 1 & 1 & 1 & 3\\
After CH02 moderation & \code{33ca044} & 1 & 1 & 1 & 3\\
After CH03 premium access & \code{7ecfa65} & 1 & 1 & 1 & 3\\
\bottomrule
\end{tabular}

\textbf{Status classification.} The LLM Provider Streaming Strategy and Auth
State Observer are preserved through these changes. The RBAC Permission Table
is also preserved as an instance: \code{rolePermissions}, \code{can},
\code{canAccess} and \code{canAny} are not replaced by a new GoF structure.
Moderation and premium access extend the application's authorization surface
through explicit checks and RLS, but are not recorded as separate RBAC-pattern
instances without evidence of a new, centralized role/permission collaboration.
The recommendation ranker likewise remains a normal helper rather than a
formal Strategy. The feature snapshots and counts must be reflected in the
final machine-readable \code{pattern_evolution.csv} and \code{Pattern_count.csv};
the current CSVs need a final consistency check against these snapshots before
submission.

\subsection*{Supporting refactoring analysis}
In addition to the three required feature changes, three refactoring experiments
were recorded on separate application branches. These are supplementary
refactoring evidence and are not the CH01--CH03 feature changes above.

\begin{tabular}{@{}p{3.1cm}p{4.9cm}p{5.9cm}@{}}
\toprule
Refactoring evidence & Code smell and technique & Pattern/verification note\\
\midrule
\code{e7dfd1c} & Repeated model-name inference; extract \code{inferProviderFromModel()} and use it in \code{generateBlogPost()}. & Streaming Strategy collaboration is preserved; this helper does not establish a new Strategy instance. The implementation record reports 23 focused tests and 1,025 tests across 88 files, while a subsequently stored focused log shows 20 tests; those records must be reconciled.\\
\code{7be15fb} & Inline title and readable-body checks in publication validation; extract \code{validatePublicationTitle()} and \code{validatePublicationContent()}. & No new GoF pattern claimed. The implementation record reports 30 focused tests before and after and 1,022 tests across 88 files.\\
\code{6ce7157} & Duplicate JSON cleanup, parsing and error handling in \code{generateBlogPost()} and \code{generateBlogPostHeadless()}; extract \code{parseGeneratedPostJson()}. & No new GoF pattern claimed. The implementation record reports 1,022 tests across 88 files; the focused log's exact branch context needs confirmation.\\
\bottomrule
\end{tabular}

The refactoring commits have separate branch histories and are not automatically
part of the feature branch. Supporting records reside in the existing
\code{changes/Change1}, \code{changes/Change2} and \code{changes/Change3}
folders. Those folder names originated with the refactoring experiments and
should be renamed or reorganized so they cannot be confused with the required
feature records CH01--CH03.

\subsection*{LOC and remaining verification}
A consistent before/after source-lines-of-code (LOC) metric has not yet been
captured using an explicit counting rule across the same paths. This report does
not invent a LOC delta. Before final submission, count production source lines
using the same tool and rule at each stated snapshot and record the commands and
outputs. Also run the focused tests for CH02 and CH03, investigate the two
reported author-byline test failures alongside CH01, run the full test suite on
the cumulative feature branch, and apply/verify both Supabase migrations against
a test database. The baseline report's lack of a local Docker/Supabase setup
explains why runtime verification was unavailable at that point; it does not
prove the new database changes work.

%------------------------------------------------------------
\section{Part 3: Cross-language comparison}

\textbf{Objective and scope.}
Part 3 compares small, equivalent implementations of two confirmed
patterns from Part 1: \textit{LLM Provider Streaming Strategy} and
\textit{Auth State Observer}. The RBAC Permission Table is intentionally
excluded from this working comparison to follow the team's request to
use two patterns. The examples isolate pattern behaviour for controlled
testing; they are not complete ports of the Next.js application and do
not call real LLM services or Supabase.

\textbf{Language scope.}
The repository contains Python and JavaScript examples by Govind and
Java and C++ examples by Soorya for both selected patterns. Python and
JavaScript share 26 saved input/output cases. Java and C++ share a
different, smaller specification with four cases per pattern. All four
languages are present, but one common four-language behavioral check
has not yet been run.

%------------------------------------------------------------
\subsection*{Selected patterns and source mapping}

\textbf{Strategy --- LLM Provider Streaming Strategy.}
Part 1 confirms this instance in
\code{code/features/ai-assistant/llmService.ts}. The concrete provider
functions are \code{streamClaude}, \code{streamOpenAI} and
\code{streamGemini}; \code{sendMessage} selects a provider from
\code{params.provider} and delegates to the corresponding function.
The original TypeScript code has an implicit function contract and a
hard-coded conditional chain rather than a separately declared strategy
interface. Govind's Python and JavaScript examples keep function-based
strategies, using \code{run} to select a provider-specific extractor for
saved sample events. For comparison, the Java version defines an
\code{LlmStrategy} interface and concrete strategy classes, while the
C++ version uses an abstract base class and concrete implementations.
The Java and C++ examples return fixed chunks. None of the four
examples calls a provider service. These are teaching adaptations,
not claims that the TypeScript code has the same class structure.

\textbf{Observer --- Auth State Observer.}
Part 1 confirms this instance in
\code{code/features/auth/context/AuthProvider.tsx},
\code{AuthContext.tsx} and
\code{code/features/auth/hooks/useSession.ts}. The original subject is
\code{supabase.auth}, supplied by the Supabase dependency;
\code{AuthProvider} registers the callback for
\code{onAuthStateChange}, updates session state and unsubscribes in
effect cleanup. Govind's Python and JavaScript examples use an
in-memory \code{AuthSubject} and an \code{AuthStateObserver} that records
user IDs. The Java and C++ examples model an event source, listeners,
event delivery and unsubscribe behaviour. None connects to Supabase.

%------------------------------------------------------------
\subsection*{Relationship to Part 2}

The uploaded Part 2 records describe three changes: CH01 related-article
recommendations, CH02 content moderation, and CH03 premium-content
access control. They also describe supplementary refactoring
experiments involving provider inference, publication validation and
generated-post JSON parsing. These feature changes do not establish a
new GoF pattern for this comparison. CH01 and the provider-inference
refactor are related to the Strategy area, but the Part 2 implementation
notes explicitly say not every provider-selection branch was refactored.

\textbf{Source-snapshot limitation.}
The copied
\code{B14/code/features/ai-assistant/llmService.ts} in the uploaded
project still contains the baseline provider-selection branches and
does not contain the helper names described in the refactoring notes.
Consequently, the snippets below compare the verified Part 1 baseline
behaviour and do not claim to benchmark a merged, post-refactoring
source snapshot. The baseline and Part 2 feature branch commits are
recorded separately in the Part 2 section and must not be conflated.

%------------------------------------------------------------
\subsection*{Common test specification}

Java and C++ use the same cases from
\code{Part3/test-cases/common_test_cases.csv}:

\begin{center}
\begin{tabular}{@{}p{1.5cm}p{4.3cm}p{8.4cm}@{}}
\toprule
Test ID & Pattern and input case & Expected behaviour \\
\midrule
S01 & Strategy: provider \code{claude}, prompt \code{hello}
& Return chunks \code{[claude:, hello]} in order. \\

S02 & Strategy: provider \code{openai}, prompt \code{hello}
& Return chunks \code{[openai:, hello]} in order. \\

S03 & Strategy: provider \code{gemini}, prompt \code{summary}
& Return chunks \code{[gemini:, summary]} in order. \\

S04 & Strategy: unknown provider
& Reject with \code{unsupported provider: unknown}. \\

O01 & Observer: sign-in event
& Active observer receives \code{signed\_in:session-1}. \\

O02 & Observer: sign-out event
& Active observer receives \code{signed\_out:none}. \\

O03 & Observer: unsubscribe then publish
& Unsubscribed observer receives no subsequent event. \\

O04 & Observer: two active observers
& Both active observers receive the same sign-in event. \\
\bottomrule
\end{tabular}
\end{center}

\textbf{Java and C++ test results.}
The recorded output in \code{Part3/evidence/test_results.txt} shows:

\begin{center}
\begin{tabular}{@{}lcc@{}}
\toprule
Pattern & Java & C++ \\
\midrule
Strategy & 4/4 passed & 4/4 passed \\
Observer & 4/4 passed & 4/4 passed \\
\bottomrule
\end{tabular}
\end{center}

The test runners are
\code{Part3/strategy/java/StrategyTest.java},
\code{Part3/strategy/cpp/test_strategy.cpp},
\code{Part3/observer/java/ObserverTest.java} and
\code{Part3/observer/cpp/test_observer.cpp}.

\textbf{Python and JavaScript test results.}
Govind's two languages use the same JSON cases in
\code{Part3/llm_provider_streaming/fixtures.json} and
\code{Part3/auth_state_observer/fixtures.json}. Each test runner checks
the expected result and that its input was not changed. A separate
process check invokes both programs for each saved input:

\begin{center}
\begin{tabular}{@{}lccc@{}}
\toprule
Pattern & Python & JavaScript & Process check \\
\midrule
Strategy & 11/11 passed & 11/11 passed & 11/11 matched \\
Observer & 15/15 passed & 15/15 passed & 15/15 matched \\
\bottomrule
\end{tabular}
\end{center}

The combined process check passed 26/26. Its script is
\code{Part3/verify_cross_language.py}; the saved output and four
language-specific logs are in \code{Part3/evidence/tests/}.
Production snippets and run commands are in \code{Part3/README.md}.
The Python/JavaScript JSON contract differs from Java/C++ direct
provider/prompt and event/session calls. The two pass counts do not
establish one shared four-language test result.

%------------------------------------------------------------
\subsection*{Measurement method}

Soorya measured the Java and C++ examples as follows:

\begin{itemize}
    \item \textbf{Source LOC:}
    Counts production files only, excluding blank lines, comments,
    preprocessor guards and test code.

    \item \textbf{Cyclomatic complexity:}
    A function-by-function manual McCabe audit uses $M=1+D$.
    Each \code{if} and loop is one decision; each non-default
    \code{switch} case is one decision. Abstract declarations and
    library internals are excluded. The audit is documented in
    \code{Part3/evidence/complexity_audit.csv} and
    \code{Part3/evidence/complexity_audit.md}; it is not the output
    of an external static-analysis tool.

    \item \textbf{Timing:}
    Compilation is performed before timing; one warm-up run is followed
    by five measured process launches on the same environment. The
    median includes process startup and the test harness but excludes
    compilation. Raw runs are in
    \code{Part3/evidence/timing_raw.csv}; toolchain versions are in
    \code{Part3/evidence/environment.txt}.
\end{itemize}

Govind used Lizard 1.24.1 for production source lines and maximum
function complexity in Python and JavaScript. His timing script ran one
warm-up and five measured whole-process launches on Windows, including
startup and JSON input/output. The raw records are in
\code{Part3/evidence/static_metrics.json} and
\code{Part3/evidence/timings/process_timings.json}. These methods and
machines differ, so the two sets of measurements are descriptive rather
than a fair language speed ranking.

%------------------------------------------------------------
\subsection*{Results}

\begin{center}
\begin{tabular}{@{}llrrr@{}}
\toprule
Pattern & Language & Production LOC & Max function CC
& Median process time (ms) \\
\midrule
Strategy & Python & 53 & 8 & 96.870 \\
Observer & Python & 74 & 21 & 100.344 \\
Strategy & JavaScript & 47 & 9 & 57.485 \\
Observer & JavaScript & 62 & 20 & 57.960 \\
Strategy & Java & 33 & 4 & 36.116 \\
Observer & Java & 24 & 2 & 42.995 \\
Strategy & C++ & 35 & 4 & 1.115 \\
Observer & C++ & 30 & 2 & 1.070 \\
\bottomrule
\end{tabular}
\end{center}

In Python and JavaScript, most of the Observer maximum function
complexity comes from validation in \code{run}, not notification itself.
For Java and C++, the maximum-complexity functions are
\code{LlmContext.select()} (Java Strategy),
\code{LlmContext::stream()} (C++ Strategy), and
\code{AuthEventSource.publish()} for Observer. Their manual decision
counts are in \code{Part3/evidence/complexity_audit.csv}.

The two pairs were measured on different machines and with different
process boundaries. These numbers cannot rank the languages. The
snippets are small, so startup and the test harness affect timing.
The C++ Observer uses non-owning pointers in this teaching example;
production C++ would require a clear lifetime and ownership policy.

%------------------------------------------------------------
\subsection*{Interpretation and limitations}

Python and JavaScript express Strategy as functions selected by a map.
Java uses an interface and concrete classes; C++ uses an abstract base
class and virtual methods. All four Observer examples hold listeners,
notify active ones and support unsubscribe. The original app instead
uses Supabase's subscription API and React Context.

Govind's four records are in \code{data/cross_language.csv}; Soorya's
four separate records remain in \code{Part3/cross_language.csv}. The
team still needs one input/output specification, a shared four-language
test run and a consistent measurement boundary before combining them.
The Part 3 LLM record is partial: the prompt and tests are saved, but
raw output and reliable model/date details were not archived. The Part 2
feature branch also needs its own remaining verification; no result
here claims to measure its post-change application code.
\section{LLM records and human verification}

\textbf{Run and model.} Part 1 pattern analysis used an LLM chat assistant on
2026-10-10. The displayed model name, exact prompts, raw outputs and the
evaluation of each analysis are recorded in \code{Part1/llm.csv},
\code{Part1/Part1_prompt.txt} and \code{Part1/llm/}.

\textbf{Claims and evaluation.} Every pattern claim was checked against the
numbered source listings in \code{Part1/evidence/}. Corrections made after
verification: the line number given for \code{getUser()} in \code{proxy.ts} was
wrong and the true line is 34; a claim that the Users page uses \code{can()} and
the server Supabase client was unsupported, since the page calls
\code{requirePermission} and \code{getAllUsers}, so it was removed; and an initial
Observer guess for the user-confirmed notification fan-out was rejected after
reading \code{route.ts}. No generated code was used in Part 1.

\textbf{Part 2 records.} The intended feature scope is preserved in
\code{Part2/Part2_prompt.txt}; source commits and test logs are referenced in
this section and in the Part 2 change records. The raw Part 2 model-response
transcripts were not separately archived as reproducible files when this report
was updated, so this report does not invent model output or treat proposed
prompts as completed model runs. The Part 2 LLM-use record and verification
status should be checked against the actual saved chat transcripts before final
submission.

\textbf{Part 3 records.} Govind's saved task prompt is in
\code{Part3/llm/govind_prompt.txt}. Source and test checks are described
in \code{Part3/llm/verification.md}. The row in
\code{Part3/llm/llm.csv} marks the full raw assistant output, displayed
model name and interaction date unavailable. It does not claim a complete
transcript or completed human review. Those fields should only be updated
from verifiable saved history.

%------------------------------------------------------------
\section{Findings, contributions and references}

\textbf{Findings.} At the baseline commit, three pattern instances were confirmed
(Strategy, Observer, RBAC), one candidate is uncertain and seven were rejected.
The application mostly relies on framework idioms and plain functions rather than
classical GoF structures. The Strategy is function-based with an \code{if}-chain
selector and no declared interface. The Observer subject comes from a dependency.
RBAC is an architectural pattern. Of the eight areas in the brief, articles,
authors, categories, comments, authentication and newsletter subscription exist;
moderation exists only as comment deletion, and recommendations and a paywall were
not found. Limitations: the app was not run (no Docker), and behavior was verified
from source and the unit-test suite only.

\textbf{Contributions.}
\begin{tabular}{@{}p{3.3cm}p{3.6cm}p{1.2cm}p{6.6cm}@{}}
\toprule
Roll number & Name & Batch & Contribution\\
\midrule
AM.SC.U4AIE23121 & Keziah Suzaine George & AIE & Part 1 pattern analysis\\
AM.SC.U4AIE23139 & Nivedita Ramesh & AIE & Part 2 change records\\
AM.SC.U4CSE23360 & Soorya Krishna & CSE & Part 3 Java and C++ examples\\
AM.SC.U4CSE23309 & B R Govind & CSE & Part 3 Python and JavaScript examples\\
\bottomrule
\end{tabular}

\textbf{References.}
\begin{enumerate}
\item E.~Gamma, R.~Helm, R.~Johnson, J.~Vlissides, \emph{Design Patterns: Elements of Reusable Object-Oriented Software}, Addison-Wesley, 1994.
\item R.~Sandhu, E.~Coyne, H.~Feinstein, C.~Youman, ``Role-Based Access Control Models,'' \emph{IEEE Computer} 29(2), 1996.
\item Next.js Blog CMS, \url{https://github.com/thepracticalengineer/nextjs-blog-cms}, commit 90ec656, MIT license.
\item Next.js documentation, \url{https://nextjs.org/docs}.
\item Supabase documentation, \url{https://supabase.com/docs}.
\item Mermaid, \url{https://mermaid.js.org}.
\end{enumerate}

\section{Team folder and datasets}
The \code{B14/} folder contains the baseline application in \code{code/},
its baseline test evidence in \code{test_result/}, pattern discovery in
\code{Part1/}, change records in \code{Part2/}, and the cross-language
examples in \code{Part3/}. The central cross-language CSV is
\code{data/cross_language.csv}; the separate Java/C++ rows are still in
\code{Part3/cross_language.csv} until one common comparison is verified.
\end{document}
